\newenvironment{newindent}[2]{\leftskip#1
}{\leftskip0cm}

\NewDocumentCommand{\art}{m o}{\noindent\textbf{Art.\ #1.%
\IfValueT{#2}{\ [#2]}}}

\NewDocumentCommand{\ust}{s m}{\IfBooleanTF{#1}{\noindent}{\\}\textbf{#2.}}

\newcommand{\pkt}[2]{\begin{newindent}{.25cm}{#1}

    \noindent\settowidth{\hangindent}{#1\ }\textbf{#1} #2

\end{newindent}}

\newcommand{\lit}[2]{\begin{newindent}{.75cm}{#1}

    \noindent\settowidth{\hangindent}{#1\ }\textbf{#1} #2

\end{newindent}}

\newcommand{\tiret}[2]{\begin{newindent}{1.25cm}{#1}

    \noindent\settowidth{\hangindent}{#1\ }\textbf{#1} #2

\end{newindent}}

%       SKŁADNIA
%
%       \art{numer}[nazwa]
%
%       \ust*{numer} noindent przy 1 
%       \ust{numer} nowa linia przy 1
%
%       \pkt{numer}{cała treść}
%
%       \lit{numer}{cała treść}
%
%       \tiret{numer}{cała treść}
%
%   zawsze przed \pkt \lit i \tiret dawaj pustą linijkę!
%       ^(.+\n)(\\pkt)
%       $1\n$2